\documentclass[11pt,reqno]{amsart}
\usepackage{amsmath,amssymb,mathtools}
\usepackage[italian]{babel}
\usepackage{microtype}
\usepackage{tikz}
\usepackage{placeins}
\usetikzlibrary{arrows.meta,positioning,calc,decorations.pathreplacing}
\usepackage[letterpaper,textwidth=6.25in,textheight=8.85in,centering]{geometry}
\usepackage[hidelinks,unicode]{hyperref}
\usepackage{booktabs,array}
\hypersetup{pdftitle={Pulizia chimica del grafene trasferito con PMMA e PPC},pdfauthor={Codex, su richiesta di Mattia Apicella},pdfsubject={Analisi della letteratura e proposta sperimentale; nessun nuovo risultato di laboratorio}}
\numberwithin{equation}{section}
\renewcommand{\abstractname}{Sommario}
\renewcommand{\refname}{Riferimenti}
\renewcommand{\keywordsname}{Parole chiave}
\renewcommand{\datename}{\textit{Data}:}
\newcommand{\SiO}{\ensuremath{\mathrm{SiO}_2}}
\newcommand{\HCl}{\ensuremath{\mathrm{HCl}}}
\newcommand{\degC}{\ensuremath{^{\circ}\mathrm{C}}}
\newcommand{\um}{\ensuremath{\mu\mathrm{m}}}
\newcommand{\nm}{\ensuremath{\mathrm{nm}}}
\title[Pulizia chimica del grafene PMMA/PPC]{Pulizia chimica del grafene trasferito con PMMA e PPC su silicio ossidato:\\ una proposta verificabile per AR-P 672.045}
\author{Nota tecnica per Mattia Apicella e Armando Serio}
\date{6 ottobre 2026}
\thanks{Analisi e testo redatti da Codex su richiesta di Mattia Apicella. Documento di proposta, non articolo sperimentale: non sono stati trattati campioni n\'e prodotti dati nuovi. Nessuna affiliazione, revisione indipendente o priorit\`a scientifica rivendicata.}
\keywords{Grafene CVD, PMMA AR-P 672.045, PPC, residui di trasferimento, pulizia chimica}
\begin{document}
\begin{abstract}
Si considera la pulizia di grafene cresciuto su rame e trasferito su
\SiO/Si con PMMA AR-P 672.045 e un ulteriore strato di PPC, dopo
un trattamento a 90\,\degC\ per due minuti e una separazione
elettrochimica dal rame. Il quesito riguarda campioni gi\`a trasferiti
e ammette solventi organici o acidi, escludendo basi e trattamenti fisici
come soluzione. La proposta principale \`e verificare un trattamento
acquoso con HCl, per il quale esiste un precedente su grafene con
residui di PMMA, mediante un confronto controllato sul materiale
effettivamente usato. Una seconda prova distingue il contributo
dell'acido da quello di una successiva estrazione in anisolo.
Si esplicitano le differenze dal precedente, il possibile ruolo del PPC,
le misure necessarie e i criteri di fallimento. Il risultato di questo
documento \`e una proposta chimica motivata e falsificabile; l'efficacia
sul campione in questione resta da dimostrare.
\end{abstract}
\maketitle

\section{La risposta al quesito}

\textbf{La prima prova che sceglierei \`e un trattamento post-trasferimento
con HCl acquoso, accompagnato da un controllo in sola acqua e da una
verifica chimica dei residui.} Non assumerei che ogni particella sia PMMA,
n\'e proporrei semplicemente di prolungare il bagno nel solvente che ha
gi\`a fallito. Il precedente pertinente \`e Xiao, Wan e Durkan
\cite{Xiao2019}; il suo trasferimento al sistema PMMA/PPC qui descritto
\`e una nostra ipotesi sperimentale.

Occorre distinguere due affermazioni. La letteratura documenta trattamenti
chimici capaci di migliorare la pulizia in condizioni definite. Non
documenta, per quanto verificato qui, una procedura universale che
garantisca la rimozione completa dal particolare campione di Armando.
Una buona proposta pu\`o contraddire l'aspettativa che si sappiano suggerire
soltanto solventi generici; per contraddire l'affermazione che sul suo
campione non funzioni nulla servono misure sul suo campione.

\subsection{Dati ricevuti e dati ancora mancanti}
Le informazioni comunicate sono: grafene CVD su lamina di Cu;
substrato finale \SiO/Si; PMMA Allresist indicato come AR672.045;
90\,\degC\ per 2 minuti; separazione dal rame in una cella con acqua
ed elettrodi; ulteriore strato di PPC in anisolo, non direttamente a
contatto con il grafene. L'identificazione delle particelle come PMMA
\`e stata formulata come probabile.

Non sono noti composizione dell'elettrolita, polarit\`a, tensione,
corrente, spessori effettivi dei due polimeri, risciacqui, eventuali altre
cotture, sequenza completa dei solventi gi\`a provati e stato di
metallizzazione. Queste assenze delimitano la proposta: non costituiscono
un'autorizzazione a ricostruire arbitrariamente il processo.

\section{Materiali e corrispondenza con il lavoro citato}

La scheda Allresist identifica AR-P 672.045 come PMMA della serie 950K,
al 4,5\% in anisolo \cite{Allresist}. Questo identifica il prodotto
di partenza, non la composizione delle particelle dopo tutti i passaggi.
La formulazione commerciale non deve essere confusa con il PPC,
poli(propilene carbonato), che \`e un secondo polimero.

Il lavoro che corrisponde meglio agli autori e alla procedura indicati
\`e \emph{Ultra-clean high-mobility graphene on technologically relevant
substrates}, di Tyagi e collaboratori, fra cui Vaidotas Mi\v{s}eikis
e Camilla Coletti \cite{Tyagi2022}. Il supplemento descrive un trasferimento
con PMMA/PPC, passaggi a 90\,\degC\ per 2 minuti e delaminazione
elettrochimica. Riporta PMMA di circa 100\,\nm\ e PPC di 1,5\,\um:
gli spessori non si possono attribuire al campione attuale. Analogamente,
l'elettrolita alcalino riportato dagli autori non \`e un dato confermato
per la cella usata da Armando.

Quel lavoro studia acetone seguito da AR 600-71. Il supplemento comprende
anche un diverso trasferimento mediante attacco del rame, con
AR-P 679.02. Non si devono scambiare i due procedimenti o i due prodotti.
Poich\'e il quesito esclude una risposta basata sul solito solvente
commerciale, quel risultato \`e qui un riferimento di confronto,
non la proposta risolutiva.

\subsection{Perch\'e il PPC cambia la diagnosi}
Essere sopra il PMMA durante il trasferimento non dimostra che il PPC
sia assente dal grafene dopo lo scioglimento dei supporti. Accesso dai
bordi, difetti della copertura e rideposizione durante il lavaggio sono
possibilit\`a da verificare. Non sono osservazioni sul campione.
Un residuo misto pu\`o richiedere una sequenza diversa da quella adeguata
al PMMA da solo. Il PPC non va inoltre confuso con il policarbonato
aromatico al bisfenolo A usato in altri metodi di trasferimento.

\section{Tre errori chimici da evitare}

\subsection{Solubilit\`a del prodotto e rimozione dell'interfaccia}
Il fatto che il resist sia venduto in anisolo dimostra che il materiale
di partenza pu\`o essere formulato in quel solvente. Non dimostra che
una frazione aderente, modificata o confinata all'interfaccia si rimuova
allo stesso modo. In un bagno reale contano anche accesso del liquido,
tempo di distacco, concentrazione del materiale estratto e
rideposizione durante lo scambio dei liquidi.

Come semplice modello diagnostico, si pu\`o scrivere
\begin{equation}\label{eq:plateau}
 m(t)=m_{\mathrm{persistente}}+
       m_{\mathrm{accessibile}}e^{-kt},\qquad k>0.
\end{equation}
Qui $m$ indica una quantit\`a di residuo, non una misura gi\`a acquisita;
``persistente'' significa tale nella finestra sperimentale. Il modello
illustra perch\'e pi\`u tempo nello stesso bagno possa produrre guadagni
sempre minori. Non identifica il meccanismo: una coda lenta, una
miscela di popolazioni o un equilibrio di riadsorbimento possono imitare
lo stesso andamento. Senza dati non si stimano $k$ o le due masse.

\subsection{Un trattamento blando non prova la carbonizzazione}
I 90\,\degC\ per 2 minuti comunicati non bastano per diagnosticare
carbonizzazione o reticolazione del PMMA. Uno studio su PMMA 950K
al 4,5\% mostra che la storia termica influisce sui residui
\cite{Jia2016}; non autorizza a trasferire al caso presente conclusioni
ottenute con cotture pi\`u severe. Servono storia completa e analisi
del residuo, non una diagnosi ricavata dalla sola insolubilit\`a apparente.

\subsection{Idrolisi dell'estere e rottura della catena sono cose diverse}
Nel PMMA gli esteri sono gruppi laterali rispetto alla catena carboniosa.
La trasformazione schematica
\begin{equation}\label{eq:hydrolysis}
 \mathrm{R{-}COOCH_3+H_2O}
 \ \rightleftharpoons\ 
 \mathrm{R{-}COOH+CH_3OH}
 \quad(\text{catalisi acida})
\end{equation}
non spezza, da sola, la catena principale. Non \`e quindi corretto
descrivere un bagno di HCl come una depolimerizzazione dimostrata del
PMMA. La formula non prova neppure che la reazione proceda in misura
utile nelle condizioni proposte. Un eventuale miglioramento potrebbe
dipendere da distacco, modifiche di bagnabilit\`a, eliminazione di
contaminanti associati o altri processi. Per il PPC la struttura della
catena \`e diversa: anche per questo occorrono riferimenti separati.

\section{Che cosa sostengono effettivamente i precedenti}

\subsection{Il precedente scelto per il campione gi\`a trasferito}
Xiao e collaboratori \cite{Xiao2019} studiano grafene su \SiO/Si
con PMMA 495K A4, diverso dall'AR-P 672.045. Dopo la rimozione
principale del supporto, impiegano HCl acquoso 1\,mol/L per 48 ore e
risciacquo in acqua deionizzata. La rugosit\`a media riportata scende
da circa 2,4 a 0,4\,\nm. Il risultato non \`e universale: su residui
invecchiati tre mesi e riscaldati a 200\,\degC\ per tre ore, la
copertura passa dal 75\% al 34\%. Il trattamento migliora quel caso
senza pulirlo completamente. Il manoscritto consultato non specifica
la temperatura del bagno. Non trasferiamo qui i risultati elettrici,
ottenuti in un diverso passaggio di fabbricazione, allo stesso campione
del confronto morfologico.

\subsection{Perch\'e non basta dire ``acido acetico''}
Liu e collaboratori \cite{Liu2023} riportano meno residui visibili
con acido acetico che con acetone, ma confermano residui dopo entrambi.
Inoltre, la tabella XPS del monostrato non sostiene una diminuzione
uniforme di tutte le componenti ossigenate: C--O passa da 7,3\% a
14,24\%, mentre O--C=O passa da 5,7\% a 6,35\%.
Il testo interpretativo e questi numeri richiedono cautela; non sono una
dimostrazione di pulizia completa. Non uso quindi quel lavoro per
promettere che un altro bagno acido risolva il campione attuale.

\subsection{Prevenzione e recupero non sono intercambiabili}
Il lavoro di Liu e collaboratori del 2025 \cite{Liu2025} esamina
contaminanti sulla faccia inferiore del film PMMA/grafene, prima della
deposizione sul substrato. I trattamenti chiamati sinteticamente
``HCl'', ``HNO$_3$'' e ``H$_2$SO$_4$'' sono miscele con perossido
di idrogeno e acqua, non i soli acidi. I risultati indicano un ruolo
dei contaminanti carboniosi gi\`a presenti durante il trasferimento.
Non provano che HCl da solo recuperi una superficie gi\`a trasferita.
La sequenza completa non viene qui adottata come procedura senza basi.

\section{Proposta per AR-P 672.045 con supporto di PPC}

\subsection{Campo di applicazione}
La prova riguarda campioni di grafene gi\`a trasferiti, con il grosso
dei supporti rimosso mediante la procedura abituale del laboratorio.
Si usano campioni gemelli, prima della metallizzazione o con la
compatibilit\`a dei metalli verificata separatamente. Non si assume
che un acido compatibile con il sistema grafene/ossido sia compatibile
con ogni contatto metallico. HF, basi, plasma, UV e ricotture non
appartengono alla proposta.

\subsection{Prima scelta: prova acida controllata}
Si confrontano un campione in HCl acquoso 1\,mol/L per 48 ore
e un gemello in acqua deionizzata per uguale durata. Per la prima
esecuzione si propone di mantenere entrambi a 20--25\,\degC,
registrando la temperatura effettiva: questa \`e una scelta del presente
progetto, non una condizione ricavata dall'articolo.
Si impiegano reagenti puliti, recipienti compatibili e le normali
procedure del laboratorio. Non si aggiungono ossidanti e non si
esegue una neutralizzazione con basi sulla superficie.

Al termine, si effettuano scambi successivi con acqua deionizzata
fresca, documentando numero, volume e durata dei bagni. L'ultimo
risciacquo si confronta con un bianco di acqua e recipiente mediante
pH e conducibilit\`a. Questo controllo limita il trascinamento del
bagno, ma non certifica l'assenza di cloro adsorbito. L'asciugatura
e il tempo prima delle misure devono essere uguali per i due campioni.

Un risultato favorevole in questa prova giustifica la ripetizione.
Un risultato negativo documenta il fallimento nelle condizioni provate,
non l'impossibilit\`a di ogni pulizia chimica. Se Armando ha gi\`a
eseguito la stessa sequenza con controlli equivalenti, non sarebbe
informativo ripeterla: si passa direttamente all'identificazione
chimica del residuo e al controllo del PPC.

\subsection{Estensione proposta: separare acido ed estrazione}
Per verificare se il trattamento renda una frazione del residuo pi\`u
estraibile, proponiamo il confronto a quattro condizioni della
Tabella~\ref{tab:design}. Questa combinazione non \`e presentata come
metodo gi\`a convalidato. L'anisolo serve come solvente definito dei
materiali iniziali e come controllo dell'estrazione; non \`e un nuovo
``solvente miracoloso''.

\begin{table}[htbp]
\caption{Confronto proposto sullo stesso processo di trasferimento.
La prima colonna identifica i gruppi, non risultati osservati.}
\label{tab:design}
\centering
\small
\begin{tabular}{@{}p{0.09\textwidth}p{0.24\textwidth}p{0.56\textwidth}@{}}
\toprule
Gruppo & Fase acquosa & Fase successiva \\
\midrule
W & Acqua deionizzata & Risciacquo, scambio comune, due bagni in isopropanolo. \\[4pt]
S & Acqua deionizzata & Risciacquo, scambio comune, due bagni in anisolo. \\[4pt]
A & HCl & Risciacquo, scambio comune, due bagni in isopropanolo. \\[4pt]
AS & HCl & Risciacquo, scambio comune, due bagni in anisolo. \\
\bottomrule
\end{tabular}
\end{table}

La fase acquosa ha uguale durata nei quattro gruppi. Dopo i risciacqui
in acqua si propone uno scambio comune in due bagni freschi di
isopropanolo, due minuti ciascuno. Segue la fase centrale: due bagni
freschi di 15 minuti ciascuno, in anisolo per S e AS, in isopropanolo
per W e A. Infine tutti ricevono due bagni freschi di isopropanolo,
due minuti ciascuno, e asciugatura con azoto. Queste durate sono
parametri iniziali proposti, non un'ottimizzazione gi\`a dimostrata.
Volumi, temperatura e modalit\`a di trasferimento vanno fissati uguali.
Si verifica su un bianco che il passaggio all'anisolo non lasci una
fase acquosa separata. W non \`e dunque un campione lasciato soltanto
in acqua: \`e il controllo dell'intera sequenza senza acido o anisolo.

Il passaggio finale a un liquido in cui il polimero \`e meno solubile
pu\`o causare rideposizione. Va quindi preceduto da ricambi sufficienti
del solvente estrattivo e verificato raccogliendo l'ultimo eluato e
controllando un campione bianco. Se il percorso di risciacquo sporca
il bianco, la sequenza AS non \`e interpretabile e va corretta prima
di usarla sul materiale di interesse.

Il confronto S contro W misura l'effetto dell'anisolo al posto
dell'isopropanolo nella fase centrale; A contro W quello acido;
AS contro entrambi stabilisce se
la sequenza combinata offra un vantaggio. Un vantaggio di AS non prova
automaticamente idrolisi o sinergia molecolare: anche la diversa
bagnabilit\`a o la rimozione di contaminanti misti sono spiegazioni
possibili.

\section{Misurare la pulizia senza confonderla con altri cambiamenti}

\subsection{Identit\`a e posizione del residuo}
Si preparano riferimenti con AR-P 672.045, PPC e i due materiali
sovrapposti, usando gli stessi lotti e la stessa storia termica. Si
includono \SiO/Si senza grafene e, se disponibile, grafene di
riferimento preparato senza contatto con questi supporti.
I riferimenti spessi servono a riconoscere le firme chimiche;
non dimostrano da soli la rimovibilit\`a di uno strato interfaciale.
Si analizzano anche riferimenti esposti agli stessi bagni: un polimero
trasformato dall'acido pu\`o perdere la firma dell'estere e rimanere
sulla superficie. La sola scomparsa di quella firma non vale come
rimozione del materiale.

La microscopia a forza atomica (AFM) localizza le particelle, ma non
ne determina la composizione. La spettroscopia fotoelettronica XPS
misura il carbonio e altre specie, ma un solo picco carbonilico non
separa PMMA e PPC. Per una distinzione credibile si confrontano
spettri completi dei riferimenti, oppure impronte di frammenti mediante
spettrometria di massa superficiale (ToF-SIMS), tenendo conto dei
frammenti comuni e degli effetti della matrice. Quando disponibile,
una misura infrarossa locale pu\`o integrare il confronto estere/carbonato.
Le analisi distruttive, come ToF-SIMS, si eseguono su provini gemelli
dedicati: non si bombarda prima l'area che si vuole poi usare per
valutare la sola pulizia chimica.

Si cercano Cu e le specie effettivamente presenti nell'elettrolita;
non si presume Fe soltanto perch\'e compare in procedimenti che
sciolgono il rame. Si controlla anche il cloro dopo il trattamento.
Il silicio osservato su \SiO\ non basta a diagnosticare contaminazione
da silicone. Attribuire una particella alla faccia inferiore del grafene
richiede una prova specifica: la sola topografia della faccia superiore
non risolve questa localizzazione.

\subsection{Quantit\`a di residuo e integrit\`a del grafene}
Prima e dopo il trattamento si acquisiscono immagini nelle stesse
aree, pi\`u aree casuali predefinite, includendo centro e bordi. Si usa
AFM con interazione debole: la scansione stessa non deve diventare
una pulizia meccanica. Un gemello misurato soltanto alla fine controlla
questo possibile artefatto.

Si riportano sia la frazione di superficie coperta da particelle sia
il volume apparente delle particelle per unit\`a di area. La sola
rugosit\`a pu\`o diminuire se il materiale si distende in un film
uniforme. Viceversa, un film uniforme pu\`o sfuggire al conteggio
delle particelle e richiede una misura chimica indipendente.
Pieghe, fori e particelle vanno distinti con criteri fissati in anticipo.

La mappa Raman controlla permanenza del grafene, segnale dei difetti
e variazioni dei picchi G e 2D nelle medesime condizioni ottiche.
L'assenza di un aumento misurabile del picco D \`e un controllo utile,
non una dimostrazione di assenza di qualunque difetto.
Si verificano inoltre delaminazioni, nuovi fori e, su substrati testimone,
eventuali modifiche dell'ossido. Misure elettriche aggiuntive devono
separare mobilit\`a, resistenza di contatto e spostamento del punto
di neutralit\`a: una variazione del drogaggio non equivale a pulizia.

\subsection{Criteri dichiarati prima della prova}
Per questa proposta si adotta una soglia operativa iniziale:
riduzione almeno del 90\% di copertura e volume apparente dei residui,
con diminuzione coerente del materiale organico, inclusi eventuali
prodotti di trasformazione, e conservazione di
almeno il 99\% dell'area iniziale di grafene. Questi numeri sono
\emph{obiettivi progettuali}, non prestazioni osservate o soglie
universali della disciplina. Il laboratorio pu\`o sostituirli prima
della prova con requisiti adeguati al dispositivo.

La riduzione dei residui si calcola sulle sole regioni con grafene
presente sia prima sia dopo, confrontando esattamente quelle regioni
anche nella misura iniziale. Le zone asportate si riportano separatamente
come perdita d'integrit\`a e non si contano come pulite. Altrimenti,
perdere una piccola regione particolarmente sporca potrebbe simulare
un grande miglioramento pur rispettando la soglia sull'area totale.

Per i difetti Raman e l'ossido si fissano margini di accettazione
in base a ripetibilit\`a strumentale e requisiti del dispositivo, prima
di vedere i risultati. Un semplice ``nessun cambiamento significativo''
in una misura poco sensibile non basta. Per dichiarare residui non
rilevati occorre riportare il limite di rivelabilit\`a, ricavato dai
bianchi e dai riferimenti, e la superficie realmente campionata.
La frase ``pulito al limite delle misure'' non significa ``privo di
qualsiasi atomo estraneo''.

\section{Esperimento minimo e decisioni successive}

Una prima esplorazione pu\`o usare quattro porzioni comparabili dello
stesso trasferimento assegnate casualmente a W, S, A e AS, con misure
iniziali e finali. Non \`e una prova di riproducibilit\`a. Se un gruppo
supera i criteri, lo si conferma con almeno tre trasferimenti
indipendenti, ciascuno suddiviso negli stessi gruppi. L'unit\`a
indipendente \`e il trasferimento, non ogni immagine AFM o ogni pixel.
Il numero definitivo di repliche va adattato alla variabilit\`a della
prova preliminare. Si conservano tutti i campioni falliti e tutte
le regioni acquisite.

\begin{figure}[htbp]
\centering
\begin{tikzpicture}[
  box/.style={draw,rounded corners=1pt,align=center,font=\small,
              text width=6.3cm,inner sep=6pt},
  flow/.style={-{Stealth[length=2mm]},thick}]
\node[box,text width=7.3cm] (q) at (0,0) {Stesso grafene e stesso processo PMMA/PPC\\
Diagnosi iniziale e confronto W, S, A, AS};
\node[box,text width=7.3cm] (m) at (0,-1.7) {Residui, firma chimica, grafene conservato\\
e danno misurati insieme};
\node[box,text width=6.6cm] (yes) at (-3.9,-3.5)
  {Criteri superati\\Ripetere su trasferimenti indipendenti};
\node[box,text width=6.6cm] (no) at (3.9,-3.5)
  {Criteri falliti\\Identificare la frazione persistente};
\draw[flow] (q)--(m);
\draw[flow] (m.south)--(yes.north);
\draw[flow] (m.south)--(no.north);
\end{tikzpicture}
\caption{Percorso decisionale proposto. Lo schema non contiene
risultati sperimentali o percentuali di successo previste.}
\end{figure}
\FloatBarrier

I confronti devono tener conto anche della quantit\`a iniziale di residuo.
Se S migliora rispetto a W senza un vantaggio aggiuntivo dell'acido,
non si attribuisce il merito a quest'ultimo. Se A migliora rispetto a W
e AS rispetto a S, il bagno acido ha un contributo misurabile,
ma il meccanismo resta da identificare. Se si riducono
le particelle senza riduzione del segnale polimerico, si sospetta
redistribuzione. Se scompare anche il grafene, il trattamento fallisce.
Se rimane una frazione carboniosa non identificabile con i riferimenti,
non si continua a chiamarla PMMA soltanto per abitudine.

Qualora prevalga contaminazione intrappolata sotto il grafene,
un trattamento dall'alto pu\`o avere accesso insufficiente. In quel
caso si apre un esperimento diverso sul trasferimento: pulizia dei
bagni, risciacqui prima della deposizione, supporto di PPC e storia
termica. Va dichiarato come prevenzione per nuovi campioni, senza
presentarlo come recupero del campione gi\`a preparato.

\section{Conclusione e limiti della risposta}

Per il sistema descritto esiste una via chimica concreta da mettere
alla prova: trattamento acido controllato, verifica della frazione
estraibile e identificazione distinta di PMMA, PPC e contaminanti
del processo. La parte originale di questa nota \`e l'adattamento
critico e il disegno di verifica; non viene rivendicata l'invenzione
del trattamento con HCl o una nuova scoperta sperimentale.

\textbf{Non \`e stato dimostrato che il campione di Armando sia pulibile
con questa procedura.} \`E stato costruito un esperimento che pu\`o
stabilirlo senza scambiare rimozione del grafene, redistribuzione
dei residui o variazione del drogaggio per successo. Se la prova
riesce e si ripete, la conclusione sar\`a limitata ai materiali e alle
condizioni misurati. Se fallisce, il fallimento dovr\`a rimanere nel
resoconto e orientare la diagnosi successiva.

\begin{thebibliography}{9}

\bibitem{Allresist}
Allresist GmbH,
\emph{AR-P 630--670: PMMA-Resists}, scheda tecnica, serie AR-P 672,
prodotto AR-P 672.045.
\href{https://www.allresist.de/wp-content/uploads/2020/03/AR-P630-670_Deutsch_Allresist_Produktinformation.pdf}{Scheda del produttore}.

\bibitem{Tyagi2022}
A. Tyagi, V. Mi\v{s}eikis, L. Martini, S. Forti, N. Mishra,
Z. M. Gebeyehu, M. A. Giambra, J. Zribi, M. Fr\'egnaux,
D. Aureau, M. Romagnoli, F. Beltram e C. Coletti,
\emph{Ultra-clean high-mobility graphene on technologically relevant substrates},
Nanoscale \textbf{14} (2022), 2167--2176.
\href{https://doi.org/10.1039/D1NR05904A}{doi:10.1039/D1NR05904A}.
\href{https://www.rsc.org/suppdata/d1/nr/d1nr05904a/d1nr05904a1.pdf}{Informazioni supplementari},
in particolare pp.~1 e 5.

\bibitem{Xiao2019}
Z. Xiao, Q. Wan e C. Durkan,
\emph{Cleaning Transferred Graphene for Optimization of Device Performance},
Advanced Materials Interfaces \textbf{6} (2019), 1801794.
\href{https://doi.org/10.1002/admi.201801794}{doi:10.1002/admi.201801794}.
\href{https://www.repository.cam.ac.uk/handle/1810/291256}{Manoscritto degli autori, Universit\`a di Cambridge}.

\bibitem{Jia2016}
K. Jia, J. Luo, R. Hu, J. Zhan, H. Cao, Y. Su, H. Zhu, L. Xie,
C. Zhao, D. Chen e T. Ye,
\emph{Evaluation of PMMA Residues as a Function of Baking Temperature and
a Graphene Heat-Free-Transfer Process to Reduce Them},
ECS Journal of Solid State Science and Technology \textbf{5} (2016), P138--P141.
\href{https://doi.org/10.1149/2.0011603jss}{doi:10.1149/2.0011603jss}.

\bibitem{Liu2023}
Z. Liu, Y. Liu, W. Zheng, Y. Ding, W. Liu, Y. Wen, H. Guo e J. Hou,
\emph{Cleanliness of transferred graphene by acetone and acid},
Frontiers in Materials \textbf{10} (2023), 1279939.
\href{https://doi.org/10.3389/fmats.2023.1279939}{doi:10.3389/fmats.2023.1279939}.

\bibitem{Liu2025}
H. Liu, Y. Li, K. Yang, L. Guo, Z. Zeng e Y. Yao,
\emph{Revealing key surface contaminants via stack-flipping strategy:
investigating rinsing protocols for clean graphene transfer and
enhanced electrical performance},
Journal of Materials Chemistry C \textbf{13} (2025), 19606--19614.
\href{https://doi.org/10.1039/D5TC02259B}{doi:10.1039/D5TC02259B}.

\end{thebibliography}
\end{document}
